\chapter{Remaining FYP Plan}
\label{ch:fyp2plan}
\section{Remaining Scope}
Integrate and evaluate the vision agents, bind real outputs to the dashboard, and prepare the final demonstration. Items listed as out of scope in Chapter~\ref{ch:intro} remain future work.

\section{Remaining Modules and Features}
\begin{itemize}
\item Integrate and evaluate Detection, Classification, Engagement and Trajectory agents; detection, classification and interaction detection with the pose judge are already implemented in the interaction CLI.
\item Evaluate interaction detection (proximity, duration, pass-by filter and pose judge) on labelled clips, including close pass-bys that are not assistance; ten local clips await labels.
\item Insight Narration agent over real analytics.
\item Real data in Dashboard, Customer Info, Employee Info, Spatial Analytics, Alerts, and System Status.
\item Implement camera handoff in the next phase: associate a person leaving one camera with a person entering another using timing, direction and appearance evidence. This is outside the mid-iteration deliverables; ambiguous associations must remain unresolved. An experimental two-camera handoff script and \texttt{/camera/*} job routes already exist.
\end{itemize}

\section{Remaining Iterations}
\begin{table}[H]
\centering
\caption{Planned FYP-2 iterations}
\begin{tabular}{P{0.15\textwidth}P{0.72\textwidth}}
\toprule
Iteration & Focus\\
\midrule
2 & Vision-component integration, real data in dashboard pages and planned cross-camera association\\
3 & Evaluation on labelled clips, tuning of thresholds and heuristics, alerts on real events\\
4 & Hardening, security, performance, documentation, and final demonstration\\
\bottomrule
\end{tabular}
\end{table}

\section{Implementation Plan}
Bind agents one at a time to the dashboard, beginning with Detection and footfall, then Classification, Engagement, and Trajectory, keeping prototype data as a fallback until each agent passes its checks. Develop and evaluate cross-camera association separately from same-camera person-distance estimation.

\section{Testing and Validation Plan}
Label a set of clips for staff/customer, interactions, and zone dwell. Report precision, recall, and error analysis per agent. Re-run Appendix test cases after each integration and record results. Check performance on CPU and record throughput.

\section{Deployment Plan}
The deliverable is a local deployment run with the documented \texttt{npm run dev:all} flow and production build commands (\texttt{npm run build}, \texttt{npm start}). A pilot with a real retail partner is future work.

\section{Tentative Timeline}
\begin{table}[H]
\centering
\caption{Tentative FYP-2 timeline (to be confirmed with the supervisor)}
\begin{tabular}{P{0.2\textwidth}P{0.67\textwidth}}
\toprule
Period & Activity\\
\midrule
Weeks 1--3 & Detection and footfall integration\\
Weeks 4--6 & Classification and engagement integration\\
Weeks 7--9 & Trajectory, narration, dashboard binding\\
Weeks 10--12 & Labelled role/interaction evaluation and tuning\\
Weeks 13--14 & Hardening, security, documentation\\
Weeks 15--16 & Final testing, demonstration, and report\\
\bottomrule
\end{tabular}
\end{table}

\section{Expected Final Deliverables}
\begin{itemize}
\item Working RetailPulse system with real vision analytics and role-aware web access.
\item Evaluation report with detector and interaction results (with and without the pose judge) and CPU throughput.
\item Source repository with documentation and run instructions.
\item Final report and demonstration.
\end{itemize}
